% Options for packages loaded elsewhere
\PassOptionsToPackage{unicode}{hyperref}
\PassOptionsToPackage{hyphens}{url}
\documentclass[
]{article}
\usepackage{xcolor}
\usepackage[margin=1in]{geometry}
\usepackage{amsmath,amssymb}
\setcounter{secnumdepth}{-\maxdimen} % remove section numbering
\usepackage{iftex}
\ifPDFTeX
  \usepackage[T1]{fontenc}
  \usepackage[utf8]{inputenc}
  \usepackage{textcomp} % provide euro and other symbols
\else % if luatex or xetex
  \usepackage{unicode-math} % this also loads fontspec
  \defaultfontfeatures{Scale=MatchLowercase}
  \defaultfontfeatures[\rmfamily]{Ligatures=TeX,Scale=1}
\fi
\usepackage{lmodern}
\ifPDFTeX\else
  % xetex/luatex font selection
\fi
% Use upquote if available, for straight quotes in verbatim environments
\IfFileExists{upquote.sty}{\usepackage{upquote}}{}
\IfFileExists{microtype.sty}{% use microtype if available
  \usepackage[]{microtype}
  \UseMicrotypeSet[protrusion]{basicmath} % disable protrusion for tt fonts
}{}
\makeatletter
\@ifundefined{KOMAClassName}{% if non-KOMA class
  \IfFileExists{parskip.sty}{%
    \usepackage{parskip}
  }{% else
    \setlength{\parindent}{0pt}
    \setlength{\parskip}{6pt plus 2pt minus 1pt}}
}{% if KOMA class
  \KOMAoptions{parskip=half}}
\makeatother
\usepackage{color}
\usepackage{fancyvrb}
\newcommand{\VerbBar}{|}
\newcommand{\VERB}{\Verb[commandchars=\\\{\}]}
\DefineVerbatimEnvironment{Highlighting}{Verbatim}{commandchars=\\\{\}}
% Add ',fontsize=\small' for more characters per line
\usepackage{framed}
\definecolor{shadecolor}{RGB}{248,248,248}
\newenvironment{Shaded}{\begin{snugshade}}{\end{snugshade}}
\newcommand{\AlertTok}[1]{\textcolor[rgb]{0.94,0.16,0.16}{#1}}
\newcommand{\AnnotationTok}[1]{\textcolor[rgb]{0.56,0.35,0.01}{\textbf{\textit{#1}}}}
\newcommand{\AttributeTok}[1]{\textcolor[rgb]{0.13,0.29,0.53}{#1}}
\newcommand{\BaseNTok}[1]{\textcolor[rgb]{0.00,0.00,0.81}{#1}}
\newcommand{\BuiltInTok}[1]{#1}
\newcommand{\CharTok}[1]{\textcolor[rgb]{0.31,0.60,0.02}{#1}}
\newcommand{\CommentTok}[1]{\textcolor[rgb]{0.56,0.35,0.01}{\textit{#1}}}
\newcommand{\CommentVarTok}[1]{\textcolor[rgb]{0.56,0.35,0.01}{\textbf{\textit{#1}}}}
\newcommand{\ConstantTok}[1]{\textcolor[rgb]{0.56,0.35,0.01}{#1}}
\newcommand{\ControlFlowTok}[1]{\textcolor[rgb]{0.13,0.29,0.53}{\textbf{#1}}}
\newcommand{\DataTypeTok}[1]{\textcolor[rgb]{0.13,0.29,0.53}{#1}}
\newcommand{\DecValTok}[1]{\textcolor[rgb]{0.00,0.00,0.81}{#1}}
\newcommand{\DocumentationTok}[1]{\textcolor[rgb]{0.56,0.35,0.01}{\textbf{\textit{#1}}}}
\newcommand{\ErrorTok}[1]{\textcolor[rgb]{0.64,0.00,0.00}{\textbf{#1}}}
\newcommand{\ExtensionTok}[1]{#1}
\newcommand{\FloatTok}[1]{\textcolor[rgb]{0.00,0.00,0.81}{#1}}
\newcommand{\FunctionTok}[1]{\textcolor[rgb]{0.13,0.29,0.53}{\textbf{#1}}}
\newcommand{\ImportTok}[1]{#1}
\newcommand{\InformationTok}[1]{\textcolor[rgb]{0.56,0.35,0.01}{\textbf{\textit{#1}}}}
\newcommand{\KeywordTok}[1]{\textcolor[rgb]{0.13,0.29,0.53}{\textbf{#1}}}
\newcommand{\NormalTok}[1]{#1}
\newcommand{\OperatorTok}[1]{\textcolor[rgb]{0.81,0.36,0.00}{\textbf{#1}}}
\newcommand{\OtherTok}[1]{\textcolor[rgb]{0.56,0.35,0.01}{#1}}
\newcommand{\PreprocessorTok}[1]{\textcolor[rgb]{0.56,0.35,0.01}{\textit{#1}}}
\newcommand{\RegionMarkerTok}[1]{#1}
\newcommand{\SpecialCharTok}[1]{\textcolor[rgb]{0.81,0.36,0.00}{\textbf{#1}}}
\newcommand{\SpecialStringTok}[1]{\textcolor[rgb]{0.31,0.60,0.02}{#1}}
\newcommand{\StringTok}[1]{\textcolor[rgb]{0.31,0.60,0.02}{#1}}
\newcommand{\VariableTok}[1]{\textcolor[rgb]{0.00,0.00,0.00}{#1}}
\newcommand{\VerbatimStringTok}[1]{\textcolor[rgb]{0.31,0.60,0.02}{#1}}
\newcommand{\WarningTok}[1]{\textcolor[rgb]{0.56,0.35,0.01}{\textbf{\textit{#1}}}}
\usepackage{graphicx}
\makeatletter
\newsavebox\pandoc@box
\newcommand*\pandocbounded[1]{% scales image to fit in text height/width
  \sbox\pandoc@box{#1}%
  \Gscale@div\@tempa{\textheight}{\dimexpr\ht\pandoc@box+\dp\pandoc@box\relax}%
  \Gscale@div\@tempb{\linewidth}{\wd\pandoc@box}%
  \ifdim\@tempb\p@<\@tempa\p@\let\@tempa\@tempb\fi% select the smaller of both
  \ifdim\@tempa\p@<\p@\scalebox{\@tempa}{\usebox\pandoc@box}%
  \else\usebox{\pandoc@box}%
  \fi%
}
% Set default figure placement to htbp
\def\fps@figure{htbp}
\makeatother
\setlength{\emergencystretch}{3em} % prevent overfull lines
\providecommand{\tightlist}{%
  \setlength{\itemsep}{0pt}\setlength{\parskip}{0pt}}
\usepackage{bookmark}
\IfFileExists{xurl.sty}{\usepackage{xurl}}{} % add URL line breaks if available
\urlstyle{same}
\hypersetup{
  pdftitle={Metropolis-Hastings},
  hidelinks,
  pdfcreator={LaTeX via pandoc}}

\title{Metropolis-Hastings}
\author{}
\date{\vspace{-2.5em}2026-04-17}

\begin{document}
\maketitle

\subsection{Introduction}\label{introduction}

The Metropolis-Hastings algorithm is the original general-purpose MCMC
sampler and remains the conceptual reference against which all
subsequent algorithms are explained. Given only an unnormalised target
density, it produces a Markov chain whose stationary distribution is
that target by repeatedly proposing a candidate value and accepting or
rejecting it based on a simple density-ratio rule. The trade-off is
universality versus efficiency: MH works for any target one can
evaluate, but tuning proposals for high-dimensional or strongly
correlated posteriors becomes painful, which is why HMC has displaced
random-walk MH for most modern Bayesian software.

\subsection{Prerequisites}\label{prerequisites}

A working understanding of Markov chains, conditional probability, and
the role of the unnormalised posterior in Bayesian inference.
Familiarity with R's distribution functions (\texttt{rnorm},
\texttt{dunif}) is needed for the example.

\subsection{Theory}\label{theory}

Given the current state \(\theta\), one iteration of MH proceeds as:

\begin{enumerate}
\def\labelenumi{\arabic{enumi}.}
\tightlist
\item
  Propose \(\theta^*\) from a proposal kernel
  \(q(\theta^* \mid \theta)\).
\item
  Compute the acceptance probability
  \(\alpha = \min\!\left(1, \dfrac{p(\theta^*) q(\theta \mid \theta^*)}{p(\theta) q(\theta^* \mid \theta)}\right)\).
\item
  With probability \(\alpha\) set the next state to \(\theta^*\);
  otherwise keep \(\theta\).
\end{enumerate}

For a symmetric proposal (Gaussian random walk centred at the current
state) the proposal terms cancel and the acceptance ratio is simply the
ratio of target densities. Detailed balance ensures the chain has the
target as its invariant distribution; ergodicity (irreducibility plus
aperiodicity) ensures convergence regardless of the starting value.

\subsection{Assumptions}\label{assumptions}

The unnormalised target density is computable everywhere in the support.
The proposal kernel must give the chain a positive probability of
reaching every part of the support; otherwise the chain is reducible and
never explores the full posterior.

\subsection{R Implementation}\label{r-implementation}

\begin{Shaded}
\begin{Highlighting}[]
\FunctionTok{set.seed}\NormalTok{(}\DecValTok{2026}\NormalTok{)}

\CommentTok{\# Target: Beta(4, 8) (unnormalised density is fine)}
\NormalTok{log\_target }\OtherTok{\textless{}{-}} \ControlFlowTok{function}\NormalTok{(x) \{}
  \ControlFlowTok{if}\NormalTok{ (x }\SpecialCharTok{\textless{}=} \DecValTok{0} \SpecialCharTok{||}\NormalTok{ x }\SpecialCharTok{\textgreater{}=} \DecValTok{1}\NormalTok{) }\FunctionTok{return}\NormalTok{(}\SpecialCharTok{{-}}\ConstantTok{Inf}\NormalTok{)}
\NormalTok{  (}\DecValTok{4} \SpecialCharTok{{-}} \DecValTok{1}\NormalTok{) }\SpecialCharTok{*} \FunctionTok{log}\NormalTok{(x) }\SpecialCharTok{+}\NormalTok{ (}\DecValTok{8} \SpecialCharTok{{-}} \DecValTok{1}\NormalTok{) }\SpecialCharTok{*} \FunctionTok{log}\NormalTok{(}\DecValTok{1} \SpecialCharTok{{-}}\NormalTok{ x)}
\NormalTok{\}}

\NormalTok{n\_iter }\OtherTok{\textless{}{-}} \DecValTok{10000}
\NormalTok{theta }\OtherTok{\textless{}{-}} \FunctionTok{numeric}\NormalTok{(n\_iter)}
\NormalTok{theta[}\DecValTok{1}\NormalTok{] }\OtherTok{\textless{}{-}} \FloatTok{0.5}
\NormalTok{sigma\_prop }\OtherTok{\textless{}{-}} \FloatTok{0.1}

\ControlFlowTok{for}\NormalTok{ (i }\ControlFlowTok{in} \DecValTok{2}\SpecialCharTok{:}\NormalTok{n\_iter) \{}
\NormalTok{  prop }\OtherTok{\textless{}{-}} \FunctionTok{rnorm}\NormalTok{(}\DecValTok{1}\NormalTok{, theta[i}\DecValTok{{-}1}\NormalTok{], sigma\_prop)}
\NormalTok{  log\_ratio }\OtherTok{\textless{}{-}} \FunctionTok{log\_target}\NormalTok{(prop) }\SpecialCharTok{{-}} \FunctionTok{log\_target}\NormalTok{(theta[i}\DecValTok{{-}1}\NormalTok{])}
  \ControlFlowTok{if}\NormalTok{ (}\FunctionTok{log}\NormalTok{(}\FunctionTok{runif}\NormalTok{(}\DecValTok{1}\NormalTok{)) }\SpecialCharTok{\textless{}}\NormalTok{ log\_ratio) theta[i] }\OtherTok{\textless{}{-}}\NormalTok{ prop}
  \ControlFlowTok{else}\NormalTok{ theta[i] }\OtherTok{\textless{}{-}}\NormalTok{ theta[i}\DecValTok{{-}1}\NormalTok{]}
\NormalTok{\}}

\CommentTok{\# Diagnostic: trace + histogram}
\FunctionTok{par}\NormalTok{(}\AttributeTok{mfrow =} \FunctionTok{c}\NormalTok{(}\DecValTok{1}\NormalTok{, }\DecValTok{2}\NormalTok{))}
\FunctionTok{plot}\NormalTok{(theta, }\AttributeTok{type =} \StringTok{"l"}\NormalTok{)}
\FunctionTok{hist}\NormalTok{(theta[}\SpecialCharTok{{-}}\NormalTok{(}\DecValTok{1}\SpecialCharTok{:}\DecValTok{1000}\NormalTok{)], }\AttributeTok{breaks =} \DecValTok{50}\NormalTok{, }\AttributeTok{freq =} \ConstantTok{FALSE}\NormalTok{, }\AttributeTok{main =} \StringTok{"Posterior samples"}\NormalTok{)}
\FunctionTok{curve}\NormalTok{(}\FunctionTok{dbeta}\NormalTok{(x, }\DecValTok{4}\NormalTok{, }\DecValTok{8}\NormalTok{), }\AttributeTok{add =} \ConstantTok{TRUE}\NormalTok{, }\AttributeTok{col =} \StringTok{"red"}\NormalTok{, }\AttributeTok{lwd =} \DecValTok{2}\NormalTok{)}
\FunctionTok{par}\NormalTok{(}\AttributeTok{mfrow =} \FunctionTok{c}\NormalTok{(}\DecValTok{1}\NormalTok{, }\DecValTok{1}\NormalTok{))}

\FunctionTok{mean}\NormalTok{(theta[}\SpecialCharTok{{-}}\NormalTok{(}\DecValTok{1}\SpecialCharTok{:}\DecValTok{1000}\NormalTok{)])     }\CommentTok{\# posterior mean}
\end{Highlighting}
\end{Shaded}

\subsection{Output \& Results}\label{output-results}

A 10,000-step chain whose post-burn-in histogram aligns with the
analytic Beta(4, 8) density. The trace plot, the histogram, and the
posterior mean (around \(4 / (4+8) = 0.33\)) confirm convergence; the
autocorrelation function tells you how many independent draws the chain
is worth.

\subsection{Interpretation}\label{interpretation}

A reporting sentence: ``A random-walk Metropolis-Hastings sampler with
proposal \(\mathcal N(\theta^{(t)}, 0.1^2)\) produced 9,000 post-burn-in
draws; the histogram tracked the analytic Beta(4, 8) target, and the
posterior mean of 0.33 matched the closed-form value to within Monte
Carlo error.'' Quote the acceptance rate so readers can judge whether
the proposal was tuned reasonably.

\subsection{Practical Tips}\label{practical-tips}

\begin{itemize}
\tightlist
\item
  Tune the proposal scale for an acceptance rate near 25--40 \% for
  one-dimensional random walks (Roberts--Gelman--Gilks rule); for higher
  dimensions, lower rates are optimal.
\item
  Discard a burn-in of 1,000--5,000 iterations and inspect the trace
  plot for visible drift before computing summaries.
\item
  When parameters are correlated, a coordinate-wise random walk has poor
  mixing; switch to block updates with a multivariate Normal proposal
  whose covariance approximates the posterior, or move to HMC.
\item
  Use \(\hat R\) across multiple chains rather than relying on a single
  trace; a single chain cannot reveal mode-trapping.
\item
  For tall posteriors with many parameters, modern HMC samplers (Stan,
  NUTS) almost always mix dramatically better than basic MH; reach for
  MH only when the target is hard to differentiate or has discrete
  components.
\end{itemize}

\subsection{R Packages Used}\label{r-packages-used}

Base R is enough for hand-coded MH; \texttt{MCMCpack} and \texttt{mcmc}
provide tuned implementations for canonical models, and \texttt{coda}
contributes diagnostic plots and convergence statistics applicable to MH
chains from any source.

\subsection{For Reviewers}\label{for-reviewers}

\textbf{What to look for in a paper using this method.}

\begin{itemize}
\tightlist
\item
  \textbf{Common misapplications.}
\item
  \textbf{Diagnostics that should be reported but often aren't.}
\item
  \textbf{Red flags in tables and figures.}
\item
  \textbf{What to verify.}
\item
  \textbf{What an adequate Methods paragraph must contain.}
\end{itemize}

\subsection{See also --- labs in this
chapter}\label{see-also-labs-in-this-chapter}

\begin{itemize}
\tightlist
\item
  \href{labs/lab-bayesian-thinking-control-vs-decision-error.Rmd}{Bayesian
  thinking; α control vs decision error}
\item
  \href{labs/lab-brms-stan-loo-hierarchical-models.Rmd}{brms / Stan,
  LOO, hierarchical models}
\end{itemize}

\end{document}
